\IfFileExists{stacks-project.cls}{%
\documentclass{stacks-project}
}{%
\documentclass{amsart}
}

% For dealing with references we use the comment environment
\usepackage{verbatim}
\newenvironment{reference}{\comment}{\endcomment}
%\newenvironment{reference}{}{}
\newenvironment{slogan}{\comment}{\endcomment}
\newenvironment{history}{\comment}{\endcomment}

% For commutative diagrams we use Xy-pic
\usepackage[all]{xy}

% We use 2cell for 2-commutative diagrams.
\xyoption{2cell}
\UseAllTwocells

% We use multicol for the list of chapters between chapters
\usepackage{multicol}

% This is generally recommended for better output
\usepackage{lmodern}
\usepackage[T1]{fontenc}

% For cross-file-references
\usepackage{xr-hyper}

% Package for hypertext links:
\usepackage{hyperref}

% For any local file, say "hello.tex" you want to link to please
% use \externaldocument[hello-]{hello}
\externaldocument[introduction-]{introduction}
\externaldocument[conventions-]{conventions}
\externaldocument[sets-]{sets}
\externaldocument[categories-]{categories}
\externaldocument[topology-]{topology}
\externaldocument[sheaves-]{sheaves}
\externaldocument[sites-]{sites}
\externaldocument[stacks-]{stacks}
\externaldocument[fields-]{fields}
\externaldocument[algebra-]{algebra}
\externaldocument[brauer-]{brauer}
\externaldocument[homology-]{homology}
\externaldocument[derived-]{derived}
\externaldocument[simplicial-]{simplicial}
\externaldocument[more-algebra-]{more-algebra}
\externaldocument[smoothing-]{smoothing}
\externaldocument[modules-]{modules}
\externaldocument[sites-modules-]{sites-modules}
\externaldocument[injectives-]{injectives}
\externaldocument[cohomology-]{cohomology}
\externaldocument[sites-cohomology-]{sites-cohomology}
\externaldocument[dga-]{dga}
\externaldocument[dpa-]{dpa}
\externaldocument[sdga-]{sdga}
\externaldocument[hypercovering-]{hypercovering}
\externaldocument[schemes-]{schemes}
\externaldocument[constructions-]{constructions}
\externaldocument[properties-]{properties}
\externaldocument[morphisms-]{morphisms}
\externaldocument[coherent-]{coherent}
\externaldocument[divisors-]{divisors}
\externaldocument[limits-]{limits}
\externaldocument[varieties-]{varieties}
\externaldocument[topologies-]{topologies}
\externaldocument[descent-]{descent}
\externaldocument[perfect-]{perfect}
\externaldocument[more-morphisms-]{more-morphisms}
\externaldocument[flat-]{flat}
\externaldocument[groupoids-]{groupoids}
\externaldocument[more-groupoids-]{more-groupoids}
\externaldocument[etale-]{etale}
\externaldocument[chow-]{chow}
\externaldocument[intersection-]{intersection}
\externaldocument[pic-]{pic}
\externaldocument[weil-]{weil}
\externaldocument[adequate-]{adequate}
\externaldocument[dualizing-]{dualizing}
\externaldocument[duality-]{duality}
\externaldocument[discriminant-]{discriminant}
\externaldocument[derham-]{derham}
\externaldocument[local-cohomology-]{local-cohomology}
\externaldocument[algebraization-]{algebraization}
\externaldocument[curves-]{curves}
\externaldocument[resolve-]{resolve}
\externaldocument[models-]{models}
\externaldocument[functors-]{functors}
\externaldocument[equiv-]{equiv}
\externaldocument[pione-]{pione}
\externaldocument[etale-cohomology-]{etale-cohomology}
\externaldocument[proetale-]{proetale}
\externaldocument[relative-cycles-]{relative-cycles}
\externaldocument[more-etale-]{more-etale}
\externaldocument[trace-]{trace}
\externaldocument[crystalline-]{crystalline}
\externaldocument[spaces-]{spaces}
\externaldocument[spaces-properties-]{spaces-properties}
\externaldocument[spaces-morphisms-]{spaces-morphisms}
\externaldocument[decent-spaces-]{decent-spaces}
\externaldocument[spaces-cohomology-]{spaces-cohomology}
\externaldocument[spaces-limits-]{spaces-limits}
\externaldocument[spaces-divisors-]{spaces-divisors}
\externaldocument[spaces-over-fields-]{spaces-over-fields}
\externaldocument[spaces-topologies-]{spaces-topologies}
\externaldocument[spaces-descent-]{spaces-descent}
\externaldocument[spaces-perfect-]{spaces-perfect}
\externaldocument[spaces-more-morphisms-]{spaces-more-morphisms}
\externaldocument[spaces-flat-]{spaces-flat}
\externaldocument[spaces-groupoids-]{spaces-groupoids}
\externaldocument[spaces-more-groupoids-]{spaces-more-groupoids}
\externaldocument[bootstrap-]{bootstrap}
\externaldocument[spaces-pushouts-]{spaces-pushouts}
\externaldocument[spaces-chow-]{spaces-chow}
\externaldocument[groupoids-quotients-]{groupoids-quotients}
\externaldocument[spaces-more-cohomology-]{spaces-more-cohomology}
\externaldocument[spaces-simplicial-]{spaces-simplicial}
\externaldocument[spaces-duality-]{spaces-duality}
\externaldocument[formal-spaces-]{formal-spaces}
\externaldocument[restricted-]{restricted}
\externaldocument[spaces-resolve-]{spaces-resolve}
\externaldocument[formal-defos-]{formal-defos}
\externaldocument[defos-]{defos}
\externaldocument[cotangent-]{cotangent}
\externaldocument[examples-defos-]{examples-defos}
\externaldocument[algebraic-]{algebraic}
\externaldocument[examples-stacks-]{examples-stacks}
\externaldocument[stacks-sheaves-]{stacks-sheaves}
\externaldocument[criteria-]{criteria}
\externaldocument[artin-]{artin}
\externaldocument[quot-]{quot}
\externaldocument[stacks-properties-]{stacks-properties}
\externaldocument[stacks-morphisms-]{stacks-morphisms}
\externaldocument[stacks-limits-]{stacks-limits}
\externaldocument[stacks-cohomology-]{stacks-cohomology}
\externaldocument[stacks-perfect-]{stacks-perfect}
\externaldocument[stacks-introduction-]{stacks-introduction}
\externaldocument[stacks-more-morphisms-]{stacks-more-morphisms}
\externaldocument[stacks-geometry-]{stacks-geometry}
\externaldocument[moduli-]{moduli}
\externaldocument[moduli-curves-]{moduli-curves}
\externaldocument[examples-]{examples}
\externaldocument[exercises-]{exercises}
\externaldocument[guide-]{guide}
\externaldocument[desirables-]{desirables}
\externaldocument[coding-]{coding}
\externaldocument[obsolete-]{obsolete}
\externaldocument[fdl-]{fdl}
\externaldocument[index-]{index}

% Theorem environments.
%
\theoremstyle{plain}
\newtheorem{theorem}[subsection]{Theorem}
\newtheorem{proposition}[subsection]{Proposition}
\newtheorem{lemma}[subsection]{Lemma}

\theoremstyle{definition}
\newtheorem{definition}[subsection]{Definition}
\newtheorem{example}[subsection]{Example}
\newtheorem{exercise}[subsection]{Exercise}
\newtheorem{situation}[subsection]{Situation}

\theoremstyle{remark}
\newtheorem{remark}[subsection]{Remark}
\newtheorem{remarks}[subsection]{Remarks}

\numberwithin{equation}{subsection}

% Macros
%
\def\lim{\mathop{\mathrm{lim}}\nolimits}
\def\colim{\mathop{\mathrm{colim}}\nolimits}
\def\Spec{\mathop{\mathrm{Spec}}}
\def\Hom{\mathop{\mathrm{Hom}}\nolimits}
\def\Ext{\mathop{\mathrm{Ext}}\nolimits}
\def\SheafHom{\mathop{\mathcal{H}\!\mathit{om}}\nolimits}
\def\SheafExt{\mathop{\mathcal{E}\!\mathit{xt}}\nolimits}
\def\Sch{\mathit{Sch}}
\def\Mor{\mathop{\mathrm{Mor}}\nolimits}
\def\Ob{\mathop{\mathrm{Ob}}\nolimits}
\def\Sh{\mathop{\mathit{Sh}}\nolimits}
\def\NL{\mathop{N\!L}\nolimits}
\def\CH{\mathop{\mathrm{CH}}\nolimits}
\def\proetale{{pro\text{-}\acute{e}tale}}
\def\etale{{\acute{e}tale}}
\def\QCoh{\mathit{QCoh}}
\def\Ker{\mathop{\mathrm{Ker}}}
\def\Im{\mathop{\mathrm{Im}}}
\def\Coker{\mathop{\mathrm{Coker}}}
\def\Coim{\mathop{\mathrm{Coim}}}

% Boxtimes
%
\DeclareMathSymbol{\boxtimes}{\mathbin}{AMSa}{"02}

%
% Macros for moduli stacks/spaces
%
\def\QCohstack{\mathcal{QC}\!\mathit{oh}}
\def\Cohstack{\mathcal{C}\!\mathit{oh}}
\def\Spacesstack{\mathcal{S}\!\mathit{paces}}
\def\Quotfunctor{\mathrm{Quot}}
\def\Hilbfunctor{\mathrm{Hilb}}
\def\Curvesstack{\mathcal{C}\!\mathit{urves}}
\def\Polarizedstack{\mathcal{P}\!\mathit{olarized}}
\def\Complexesstack{\mathcal{C}\!\mathit{omplexes}}
% \Pic is the operator that assigns to X its picard group, usage \Pic(X)
% \Picardstack_{X/B} denotes the Picard stack of X over B
% \Picardfunctor_{X/B} denotes the Picard functor of X over B
\def\Pic{\mathop{\mathrm{Pic}}\nolimits}
\def\Picardstack{\mathcal{P}\!\mathit{ic}}
\def\Picardfunctor{\mathrm{Pic}}
\def\Deformationcategory{\mathcal{D}\!\mathit{ef}}
\usepackage{cleveref}
\crefname{originaltheorem}{original theorem}{original theorems}
\Crefname{originaltheorem}{Original theorem}{Original theorems}
\makeatletter
\newlabel{theorem-C}{{03V2}{source}{}{sourcepointer-1}{}}
\newlabel{theorem-C@cref}{{[originaltheorem][0][]03V2}{[1][0][]source}{}{}{}}
\makeatother
\title{Compatible systems of perfect complexes and endomorphisms assemble into a limit complex with compatible traces}
\author{The Stacks Project Authors}
\date{stacks-project:trace.tex:local-snapshot-20261012; GNU FDL 1.2 or later}
\begin{document}
\maketitle
\setcounter{subsection}{1}
\noindent\textit{Editorial note.} All delivered mathematical text is the author's own, copied from the pinned Stacks Project chapter file trace.tex; no statement or proof step was rewritten, reordered or removed. The counted span is the author environment and its own proof environment. References inside the span are handed to the assembler as navigation stubs that carry the canonical tag of the cited tag only; whether each of them is load-bearing is recorded on the stub itself.\par\medskip
\noindent\textit{Allowed standard prerequisites (editorial).}\begin{itemize}
\item A perfect complex over the local ring Z/l\textasciicircum{}nZ is represented by a finite complex of finite free modules with differentials landing in l times the next term, unique up to isomorphism. invoked at proof lines 2613-2618. Provides the canonical representatives K\_n\textasciicircum{}bullet and their uniform boundedness.
\item The trace of an endomorphism of a bounded complex over Q\_l computes as the alternating sum of traces on cohomology. invoked at proof lines 2680-2687. Identifies the limit trace with the trace on H\textasciicircum{}i(K\_infinity tensor Q\_l).
\end{itemize}\medskip
\ifdefined\hypertarget\hypertarget{sourcepointer-1}{}\fi
\noindent\textit{Source reference.} Original theorem 03V2: sources/stacks-trace/pinned-original-source.tex, source lines 2523--2523. Cited at proof lines 2686-2688 to record that this computation also finishes the proof of theorem-C; navigational reference at the close of the span.\par\medskip
\begin{lemma}
\label{lemma-piece-together}
Suppose we have
$K_n\in D_{perf}(\mathbf{Z}/\ell^n\mathbf{Z})$, $\pi_n : K_n\to K_n$
and isomorphisms
$\varphi_n :
K_{n+1} \otimes^\mathbf{L}_{\mathbf{Z}/\ell^{n+1}\mathbf{Z}}
\mathbf{Z}/\ell^n\mathbf{Z}
\to K_n$
compatible with $\pi_{n+1}$ and $\pi_n$. Then
\begin{enumerate}
\item
the elements $t_n = \text{Tr}(\pi_n |_{K_n})\in \mathbf{Z}/\ell^n\mathbf{Z}$
form an element $t_\infty = \{t_n\}$ of $\mathbf{Z}_\ell$,
\item
the $\mathbf{Z}_\ell$-module $H_\infty^i = \lim_n H^i(k_n)$ is finite and
is nonzero for finitely many $i$ only, and
\item
the operators $H^i(\pi_n): H^i(K_n)\to H^i(K_n)$ are compatible and define
$\pi_\infty^i : H_\infty^i\to H_\infty^i$ satisfying
$$
\sum (-1)^i \text{Tr}(
\pi_\infty^i |_{H_\infty^i \otimes_{\mathbf{Z}_\ell}\mathbf{Q}_\ell}) =
t_\infty.
$$
\end{enumerate}
\end{lemma}

\begin{proof}
Since $\mathbf{Z}/\ell^n\mathbf{Z}$ is a local ring and $K_n$ is perfect, each
$K_n$ can be represented by a finite complex $K_n^\bullet$ of finite free
$\mathbf{Z}/\ell^n \mathbf{Z}$-modules such that the map $K_n^p \to K_n^{p+1}$
has image contained in $\ell K_n^{p+1}$. It is a fact that such a complex is
unique up to isomorphism. Moreover $\pi_n$ can be represented by a morphism of
complexes $\pi_n^\bullet : K_n^\bullet\to K_n^\bullet$ (which is unique up to
homotopy). By the same token the isomorphism
$\varphi_n : K_{n+1} \otimes_{\mathbf{Z}/\ell^{n+1}\mathbf{Z}}^{\mathbf{L}}
\mathbf{Z}/\ell^n\mathbf{Z}\to K_n$ is represented by a map of complexes
$$
\varphi_n^\bullet :
K_{n+1}^\bullet
\otimes_{\mathbf{Z}/\ell^{n+1}\mathbf{Z}}
\mathbf{Z}/\ell^n\mathbf{Z} \to K_n^\bullet.
$$
In fact, $\varphi_n^\bullet$ is an isomorphism of complexes, thus we see that
\begin{itemize}
\item
there exist $a, b\in \mathbf{Z}$ independent of $n$ such that $K_n^i = 0$ for
all $i\notin[a, b]$, and
\item
the rank of $K_n^i$ is independent of $n$.
\end{itemize}
Therefore, the module $K_\infty^i = \lim_n \{K_n^i, \varphi_n^i\}$ is a
finite free $\mathbf{Z}_\ell$-module and $K_\infty^\bullet$ is a finite complex
of finite free $\mathbf{Z}_\ell$-modules. By induction on the number of nonzero
terms, one can prove that $H^i\left(K_\infty^\bullet\right) = \lim_n
H^i\left(K_n^\bullet\right)$ (this is not true for unbounded complexes). We
conclude that $H_\infty^i = H^i\left(K_\infty^\bullet\right)$ is a finite
$\mathbf{Z}_\ell$-module. This proves {\it ii}. To prove the remainder of the
lemma, we need to overcome the possible noncommutativity of the diagrams
$$
\xymatrix{
{K_{n+1}^\bullet} \ar[d]_{\pi_{n+1}^\bullet} \ar[r]^{\varphi_n^\bullet} &
{K_n^\bullet} \ar[d]^{\pi_n^\bullet} \\
{K_{n+1}^\bullet} \ar[r]_{\varphi_n^\bullet} & {K_n^\bullet.}
}
$$
However, this diagram does commute in the derived category, hence it commutes
up to homotopy. We inductively replace $\pi_n^\bullet$ for $n\geq 2$ by
homotopic maps of complexes making these diagrams commute. Namely, if $h^i :
K_{n+1}^i \to K_n^{i-1}$ is a homotopy, i.e.,
$$
\pi_n^\bullet \circ \varphi_n^\bullet -
\varphi_n^\bullet \circ \pi_{n + 1}^\bullet = dh + hd,
$$
then we choose $\tilde h^i : K_{n+1}^i\to K_{n+1}^{i-1}$ lifting $h^i$. This is
possible because $K_{n+1}^i$ free and $K_{n+1}^{i-1}\to K_n^{i-1}$ is
surjective. Then replace $\pi_n^\bullet$ by $\tilde\pi_n^\bullet$ defined by
$$
\tilde\pi_{n+1}^\bullet = \pi_{n+1}^\bullet + d\tilde h+\tilde hd.
$$
With this choice of $\{\pi_n^\bullet\}$, the above diagrams commute, and the
maps fit together to define an endomorphism $\pi_\infty^\bullet =
\lim_n\pi_n^\bullet$ of $K_\infty^\bullet$. Then part {\it i} is clear:
the elements $t_n = \sum(-1)^i \text{Tr}\left(\pi_n^i |_{K_n^i}\right)$
fit into an element $t_\infty$ of $\mathbf{Z}_\ell$. Moreover
\begin{align*}
t_\infty
& =
\sum (-1)^i \text{Tr}_{\mathbf{Z}_\ell}(\pi_\infty^i |_{K_\infty^i}) \\
& =
\sum (-1)^i \text{Tr}_{\mathbf{Q}_\ell}(
\pi_\infty^i |_{K_\infty^i \otimes_{\mathbf{Z}_\ell}\mathbf{Q}_\ell}) \\
& =
\sum (-1)^i \text{Tr}(
\pi_\infty |_{H^i(K_\infty^\bullet \otimes \mathbf{Q}_\ell)})
\end{align*}
where the last equality follows from the fact that $\mathbf{Q}_\ell$ is a
field, so the complex $K_\infty^\bullet \otimes \mathbf{Q}_\ell$ is
quasi-isomorphic to its cohomology
$H^i(K_\infty^\bullet \otimes \mathbf{Q}_\ell)$. The latter is also equal to
$H^i(K_\infty^\bullet)\otimes_{\mathbf{Z}}\mathbf{Q}_\ell = H_\infty^i \otimes
\mathbf{Q}_\ell$, which finishes the proof of the lemma, and also that of
Theorem \ref{theorem-C}.
\end{proof}

\end{document}
